
We measure an oracle gap of 15.09\% relative improvement (11.62 percentage points absolute) between per-task optimal LoRA adapter selection and the best fixed-rank baseline across multi-domain NLP benchmarks. Oracle ranks distribute as 5 tasks selecting rank-4, 4 tasks selecting rank-8, 4 tasks selecting rank-16, and 4 tasks selecting rank-32. The best fixed rank (rank-8) achieves 76.97\% average accuracy, while the per-task oracle achieves 88.58\%.

Current practice in parameter-efficient fine-tuning applies a single LoRA rank globally across all downstream tasks. Practitioners select rank 8 or 16 based on validation performance, then apply this configuration uniformly. This approach assumes that optimal adapter capacity generalizes across heterogeneous task distributions. Our systematic measurement tests this assumption by training all combinations of 17 tasks and 4 adapter ranks, then comparing per-task optimal selection against fixed-rank baselines.

The research question is: how much performance is lost by forcing a single fixed adapter rank across heterogeneous task distributions? We address this through exhaustive evaluation on GLUE and XTREME benchmarks, measuring the oracle gap under perfect hindsight selection.

LoRA (Hu et al., 2021) injects low-rank trainable matrices into frozen transformer layers, with adapter rank controlling the capacity-efficiency trade-off. However, existing work treats rank as a global hyperparameter chosen once and applied uniformly. Prior work on Neural Architecture Search (Zoph and Le, 2017) and hyperparameter optimization (Feurer et al., 2015) has measured per-task configuration benefits in architecture and hyperparameter spaces, but no systematic measurement exists for LoRA rank oracle gap on multi-domain benchmarks.

Our hypothesis is that task heterogeneity in multi-domain benchmarks creates measurable optimization opportunities through task-specific adapter rank selection. To test this, we train all rank-task configurations (68 experiments), compute the per-task best (oracle), and compare to the best single fixed rank. The oracle gap establishes an upper bound for what task-aware routing mechanisms could achieve under perfect selection.

Our contributions are:

\textbf{First}, we provide the first systematic measurement of LoRA rank oracle gap on multi-domain NLP benchmarks, quantifying a 15.09\% relative performance improvement (11.62 percentage points absolute) between per-task optimal selection and the best fixed-rank baseline. This measurement establishes an upper bound that practical routing mechanisms cannot exceed.

\textbf{Second}, we demonstrate that oracle selections distribute across adapter ranks with a 5/4/4/4 split across ranks 4-8-16-32. This distribution indicates that different tasks select different capacity levels, rather than clustering around a single optimal rank. The best fixed rank (rank-8 at 76.97\%) underperforms the per-task oracle (88.58\%), confirming that no single configuration achieves optimal performance across heterogeneous tasks.

\textbf{Third}, we document that rank-32 achieves the worst average performance (62.95\%) despite having 8× more parameters than rank-4. Rank-32 collapses to random baseline (50\% accuracy) on CoLA (8,551 samples), SST-2 (67,349 samples), and WNLI (635 samples), demonstrating that higher adapter capacity does not guarantee better performance when the number of parameters exceeds data complexity.

The 15.09\% relative improvement represents an upper bound under perfect hindsight selection. Practical routing mechanisms would face classifier errors and computational overhead, reducing achievable improvement. Our work establishes that the oracle gap exists and quantifies its magnitude; whether routing systems can capture a substantial fraction of this gap remains an open question.

The paper proceeds as follows: Section 2 reviews parameter-efficient fine-tuning literature; Section 3 describes the oracle gap measurement methodology; Section 4 presents experimental setup; Section 5 reports results including the 15.09\% gap, oracle selection distribution, and rank-32 performance analysis; Section 6 discusses implications and limitations; Section 7 concludes.
